\clearpage
\begin{trapbox}
\begin{minipage}{\linewidth}
\paragraph{易错一：以为严格递增就必须处处 $f'(x)>0$}
\begin{center}
\begin{tikzpicture}
\begin{axis}[at={(0,0)},anchor=south west,width=6.1cm,height=5.3cm,
 axis lines=middle,xmin=-1.4,xmax=1.4,ymin=-2.3,ymax=2.3,
 ticks=none,xlabel={$x$},ylabel={$y$},title={$x^3$：零点处仍严格递增},title style={font=\small}]
\addplot[blue,very thick,domain=-1.2:1.2,samples=80]{x^3};
\addplot[orange!85!black,dashed,thick,domain=-0.7:0.7]{0};
\node[font=\small,anchor=north west] at (axis cs:0,-0.2) {$g'(0)=0$};
\end{axis}
\begin{axis}[at={(6.2cm,0)},anchor=south west,width=6.1cm,height=5.3cm,
 axis lines=middle,xmin=-1.4,xmax=1.4,ymin=-2.3,ymax=2.3,
 ticks=none,xlabel={$x$},ylabel={$y$},title={$1$：不减却不严格递增},title style={font=\small}]
\addplot[blue,very thick,domain=-1.2:1.2]{1};
\node[font=\small,anchor=west] at (axis cs:0.2,-1) {$f'\equiv0$};
\end{axis}
\end{tikzpicture}
\end{center}

\end{minipage}

\textbf{错误想法：}只要有一个位置导数为 $0$，函数就不严格递增。

\textbf{反例：}$f(x)=x^3$ 在整个实数范围内严格递增，但 $f'(0)=0$。

\textbf{正确理解：}$f'(x)>0$ 处处成立，一定严格递增；但严格递增允许某些位置导数为 $0$。如果只知道 $f'(x)\ge0$，还要防止曲线出现一整段水平线。例如常函数 $f(x)=2$ 的导数处处为 $0$，却不严格递增。

\begin{minipage}{\linewidth}
\paragraph{易错二：把被断开的两段当成一整段}
\begin{center}
\begin{tikzpicture}
\begin{axis}[width=11.6cm,height=6.2cm,axis lines=middle,
 xmin=-3,xmax=3,ymin=-3,ymax=3,xlabel={$x$},ylabel={$y$},
 xtick={-1,1},ytick={-1,1},title={$y=-1/x$：每支递增，跨过缺口却不递增},title style={font=\small}]
\addplot[blue,very thick,domain=-2.8:-0.36,samples=100]{-1/x};
\addplot[blue,very thick,domain=0.36:2.8,samples=100]{-1/x};
\addplot[only marks,orange!85!black] coordinates {(-1,1) (1,-1)};
\node[font=\small,above left] at (axis cs:-1,1) {$(-1,1)$};
\node[font=\small,below right] at (axis cs:1,-1) {$(1,-1)$};
\draw[orange!85!black,dashed,-Stealth] (axis cs:-1,1)--(axis cs:1,-1);
\end{axis}
\end{tikzpicture}
\end{center}

\end{minipage}

看 $h(x)=-\dfrac{1}{x}$，它只能取 $x\ne0$，而且
\[
h'(x)=\frac1{x^2}>0\quad(x\ne0).
\]
于是它在 $(-\infty,0)$ 内递增，在 $(0,+\infty)$ 内也递增，\textbf{却不能说在所有非零实数构成的范围内整体递增}。

直接比较两个点就明白了：
\[
-1<1,\qquad h(-1)=1>h(1)=-1.
\]
\textbf{记住：}分母为 $0$、无意义的位置，相当于路断开了；导数判定要按连续的区间分别进行。

\newpage
\begin{minipage}{\linewidth}
\paragraph{易错三：忽略端点处的函数值}
\begin{center}
\begin{tikzpicture}
\begin{axis}[width=10.8cm,height=5.5cm,axis lines=middle,
 xmin=-0.3,xmax=1.4,ymin=-0.3,ymax=2.5,
 xlabel={$x$},ylabel={$y$},xtick={1},ytick={1,2},
 title={内部斜率为正，端点跳跃仍可破坏递增},title style={font=\small}]
\addplot[blue,very thick,domain=0:1]{x};
\addplot[only marks,mark=*,mark options={fill=white,draw=blue},mark size=2.5pt] coordinates {(0,0)};
\addplot[only marks,mark=*,blue,mark size=2.5pt] coordinates {(0,2) (1,1)};
\node[font=\small,right] at (axis cs:0,2) {$f(0)=2$};
\node[font=\small,above left] at (axis cs:1,1) {$f(1)=1$};
\end{axis}
\end{tikzpicture}
\end{center}

\end{minipage}

设
\[
f(x)=\begin{cases}
2,&x=0,\\
x,&0<x\le1.
\end{cases}
\]
虽然在 $0<x<1$ 时，$f'(x)=1>0$，但是
\[
0<1,\qquad f(0)=2>f(1)=1.
\]
因此它\textbf{不能说在 $[0,1]$ 上递增}。

\textbf{原因：}在 $x=0$ 处，函数值突然跳到了 $2$，破坏了端点与内部的大小顺序。要用前述导数判据直接判断闭区间的单调性，应核对端点连续；若端点不连续，就要另行比较端点与内部函数值。\textbf{不能反过来说严格递增一定要求端点连续。}

\par\medskip
\noindent\begin{minipage}{\linewidth}
\paragraph{易错四：导数等于 $0$，就认定有最高点或最低点}
\begin{center}\begin{tikzpicture}[font=\small,>=stealth,line join=round]
\draw[gray!45,rounded corners=3pt] (0,0) rectangle (14.2,8.8);
\node at (7.1,8.2) {同样导数为零，三种图像给出不同结论};

\begin{scope}[shift={(2.5,4.6)},scale=1.05]
\draw[->,gray] (-1.4,0)--(1.45,0) node[below] {$x$};
\draw[->,gray] (0,-1.9)--(0,2) node[left] {$y$};
\draw[blue!75!black,very thick,samples=90,domain=-1.17:1.17] plot (\x,{(\x)*(\x)*(\x)});
\draw[orange!85!black,dashed,thick] (-.65,0)--(.65,0);
\fill (0,0) circle(2pt);
\end{scope}
\node at (2.5,7.3) {$f(x)=x^3$};
\node at (2.5,1.95) {$f':\ +\quad0\quad+$};
\node at (2.5,1.25) {没有极值};

\begin{scope}[shift={(7.1,4.6)},scale=1.05]
\draw[->,gray] (-1.4,0)--(1.45,0) node[below] {$x$};
\draw[->,gray] (0,-1.9)--(0,2) node[left] {$y$};
\draw[blue!75!black,very thick,samples=90,domain=-1.17:1.17] plot (\x,{-1*(\x)*(\x)});
\draw[orange!85!black,dashed,thick] (-.65,0)--(.65,0);
\fill (0,0) circle(2pt);
\end{scope}
\node at (7.1,7.3) {$f(x)=-x^2$};
\node at (7.1,1.95) {$f':\ +\quad0\quad-$};
\node at (7.1,1.25) {极大值};

\begin{scope}[shift={(11.7,4.6)},scale=1.05]
\draw[->,gray] (-1.4,0)--(1.45,0) node[below] {$x$};
\draw[->,gray] (0,-1.9)--(0,2) node[left] {$y$};
\draw[blue!75!black,very thick,samples=90,domain=-1.17:1.17] plot (\x,{(\x)*(\x)});
\draw[orange!85!black,dashed,thick] (-.65,0)--(.65,0);
\fill (0,0) circle(2pt);
\end{scope}
\node at (11.7,7.3) {$f(x)=x^2$};
\node at (11.7,1.95) {$f':\ -\quad0\quad+$};
\node at (11.7,1.25) {极小值};
\node[text width=13.2cm,align=center] at (7.1,.55) {图中三点导数均为 $0$；左图不转向，后两图分别形成峰与谷。};
\end{tikzpicture}
\end{center}
\end{minipage}\par\medskip
图中 $x^3$ 在 $0$ 处导数为零，却始终递增；$-x^2$ 与 $x^2$ 则分别在 $0$ 处形成极大值、极小值。\textbf{单凭零导数，不能区分这三种情形。}

\textbf{正确做法：}找出候选点后，先确认函数在该点连续、左右邻域可导。若左右两侧的导数各自保持固定的正负号，则左正右负可判断为\textbf{极大值}，左负右正可判断为\textbf{极小值}；两侧同号就不能用它判断为极值。这里说的是附近的极值，不是整个定义域的最大值或最小值。

\par\medskip
\noindent\begin{minipage}{\linewidth}
\paragraph{易错五：$f'(x)\ge0$ 就一定严格递增}
\begin{center}\begin{tikzpicture}[font=\small,>=stealth,line join=round]
\draw[gray!45,rounded corners=3pt] (0,0) rectangle (14.2,7.6);
\node at (7.1,7.0) {导数非负只能先说明不下降};

\begin{scope}[shift={(3.8,2.4)},scale=1.1]
\draw[->,gray] (-1.8,0)--(2.3,0) node[below] {$x$};
\draw[->,gray] (0,-.45)--(0,3.1) node[left] {$y$};
\draw[blue!75!black,very thick] (-1.6,2)--(1.9,2);
\draw[gray,dashed] (-1,0)--(-1,2) (1,0)--(1,2);
\fill (-1,2) circle(2pt) node[above left=7pt] {$(-1,2)$};
\fill (1,2) circle(2pt) node[above right=7pt] {$(1,2)$};
\node[below=6pt] at (-1,0) {$-1$};\node[below=6pt] at (1,0) {$1$};
\end{scope}
\node[anchor=west,align=left] at (8.4,4.2)
{$f(x)=2,\quad f'(x)=0$\\[16pt]
$-1<1$，但 $f(-1)=f(1)$\\[16pt]
不满足严格递增所需的\\[6pt]$f(-1)<f(1)$。};
\node[text width=13.2cm,align=center] at (7.1,.65) {常函数的两点等高；它不会下降，却不是严格递增。};

\end{tikzpicture}
\end{center}
\end{minipage}\par\medskip
图中 $f(x)=2$ 满足 $f'(x)=0\ge0$，但两个不同位置的函数值相等，不满足严格递增的定义。

\textbf{正确做法：}在前述区间的连续、可导条件下，$f'(x)\ge0$ 先说明``不会下降''。若还没有任何一整段保持不变，才能进一步得到严格递增。\textbf{特别地：全区间导数非负，且零导数点只有有限个，就能保证严格递增。}


\par\medskip
\noindent\begin{minipage}{\linewidth}
\paragraph{易错六：认为某点导数不存在，函数就不能严格递增}
\begin{center}\begin{tikzpicture}[font=\small,>=stealth,line join=round]
\draw[gray!45,rounded corners=3pt] (0,0) rectangle (14.2,8.8);
\node at (7.1,8.2) {竖直切线处没有有限导数，仍可严格递增};

\begin{scope}[shift={(3.8,4.4)},scale=1.4]
\draw[->,gray] (-2.1,0)--(2.1,0) node[below] {$x$};
\draw[->,gray] (0,-1.65)--(0,1.65) node[left] {$y$};
\draw[blue!75!black,very thick,samples=140,domain=-1.25:1.25,variable=\t] plot ({(\t)*(\t)*(\t)},\t);
\draw[orange!85!black,dashed,thick] (0,-1.4)--(0,1.4);
\draw[green!45!black,thick] (-.4,-.4)--(1.25,1.25);
\draw[purple!80!black,thick] (-.32,-1.28)--(.32,1.28);
\fill (0,0) circle(1.8pt);\node at (-.35,.2) {$O$};
\fill[green!45!black] (1,1) circle(2pt);
\fill[purple!80!black] (.125,.5) circle(2pt);
\end{scope}
\node[anchor=west,align=left] at (8.5,4.8)
{$f(x)=\sqrt[3]{x}$\\[15pt]
$\dfrac{f(h)-f(0)}h=\dfrac1{(\sqrt[3]h)^2}$\\[17pt]
$\textcolor{green!45!black}{h=1:\quad\text{差商}=1}$\\[12pt]
$\textcolor{purple!80!black}{h=\tfrac18:\quad\text{差商}=4}$\\[12pt]
$h=\tfrac1{64}:\quad\text{差商}=16$};
\node[text width=13.2cm,align=center] at (7.1,1.2) {绿、紫两条割线越来越陡；从左右趋近 $0$ 时，差商都没有有限极限。};
\node[text width=13.2cm,align=center] at (7.1,.55) {橙色虚线为竖直切线；$f'(0)$ 不存在，不能写成一个有限数。};

\end{tikzpicture}
\end{center}
\end{minipage}\par\medskip

\textbf{错误想法：}要在整段上严格递增，导数就必须在每一个位置都存在。

\textbf{反例：}考虑
\[
f(x)=\sqrt[3]{x}.
\]
它在整个实数范围内严格递增。因为任取 $a<b$，都有
\[
\sqrt[3]{a}<\sqrt[3]{b}.
\]
理由可用前面已证明的三次函数严格递增：若 $\sqrt[3]a\ge\sqrt[3]b$，取三次方便得到 $a\ge b$，与 $a<b$ 矛盾。

但在 $x=0$ 处，用导数的定义计算：
\[
\frac{f(0+h)-f(0)}{h}
=\frac{\sqrt[3]{h}}{h}
=\frac{1}{\bigl(\sqrt[3]{h}\bigr)^2}\qquad(h\ne0).
\]
当 $h$ 越来越接近 $0$ 时，右边会无限增大，不会靠近某个确定的有限数，所以 $x=0$ 处\textbf{导数不存在}。

\textbf{正确理解：}在前述连续区间和端点条件下，内部导数处处大于 $0$，足以保证函数严格递增；但要证明一个函数严格递增，\textbf{并不必须要求它处处有导数}。如果导数在个别点不存在，可以尝试直接比较两个函数值，或结合图像分段判断，不能机械地套用导数规则。
\end{trapbox}
